\section{The hardware state in a tactile image}
\label{sec:signature}

\paragraph{Data.}
A GelSight Mini (1640$\times$1232 native frames, \SI{11.3}{\micro\metre}/px, \SI{18.8}{fps}, printed markers) is mounted on one finger of a parallel gripper, and an ATI Nano17 on the other finger measures the normal force. The object is a box whose face carries one of 20 materials. It is clamped in a fixture with foam behind it so that the contact aligns itself. Each episode grasps the face at one of two forces (40 or 80\pct\ of a per-material maximum, 0.9--5\,N) and strokes it by 2\,cm in four directions and a circle. The gel is lubricated every five episodes, and one session records one material. Four campaigns cover four months: A (June 15--22), B (July 6--8), C (July 9--13) and D (September 22 to October 15), \EpRecorded\ episodes in all. \EpNoCrops\ episodes of campaign D have no eight-frame span of contact inside the stroke, so the material experiments use \EpCrops. Two units and three gels were used. Unit 1 carried gel 1 until July 9, unit 2 carried gel 2 from July 10, and gel 3 replaced gel 2 on October 9 between 18:42 and 19:37. We call these three combinations hardware states S1, S2 and S3. The laboratory log recorded a change of unit and four gel replacements. The images show exactly the two changes above, and the operators confirmed them against their notes. Campaign D also records frames of the bare gel. The camera's exposure and white balance are fixed: on a held contact the mean of each colour channel drifts by at most 0.3 grey levels per episode. We never resample a frame. Material models read native 192-px crops.

\paragraph{Measurement.}
For every session we form the position-wise mean of its contact crops in frame coordinates, its \emph{state field} $F_s(p)$. For a sampled contact crop we correlate its fine band (190--370\,\si{\micro\metre}) with the same band of another session's field, excluding pixels within 60\,px of a printed marker so that inpainting residue cannot contribute. \Cref{fig:signature} shows the medians over sessions. A contact matches its own session's other contacts at \HwFarOwn, another object touched the same day at \HwFarSameDay, the same unit and gel on another day at \HwFarWeeks, and the same unit and gel three months later at \HwFarMonths. It does not match the other unit (\HwFarOtherUnit) or the other gel of the same unit (\HwFarOtherGel). Frames of the bare gel, with nothing touching it, match their session's contact field at \BlankVsFieldFar. Another lab's GelSight Mini shows the same ordering~\cite{tactilemnist2025}, although its correlation across a gel swap stays higher (\TmSigOtherGel).

\paragraph{What the state is.}
The pattern appears to be fixed relative to the camera. Warping a bare-gel frame by the membrane's measured displacement lowers its match to the session field from \BlankVsFieldFar\ to \WarpedBlankVsFieldFar. The membrane moves by a median of only \WarpDispPx\,px, however, so this test is weak. It is nevertheless carried by the gel cartridge, because a new gel on the same unit and camera removes it. We therefore read it as an optical property of the cartridge seen through the camera, rather than as texture on the deforming surface. The state also has a smooth part, the illumination field, which changes a little from session to session on the same hardware. Between two bare-gel frames of one session, the illumination moves by a median \DriftUse\ grey levels over \DriftUseContacts\ grasps. Between the last bare-gel frame of one session and the first of the next, with no contact in between, it moves by \DriftIdle. A contact matches its own session (\HwFarOwn) better than another object touched the same day (\HwFarSameDay) because the session's mean contains the same object. Within one unit and gel, a session's illumination field lies \LpDistWithinSthree--\LpDistWithinSone\ grey levels from the mean of its hardware state. The two units differ by \LpDistSoneStwo\ grey levels, and the two gels on unit 2 by \LpDistStwoSthree.

\paragraph{The images record the hardware history.}
The changes are sharp enough to date from the images alone. Correlating every pair of the 35 June and July sessions splits them into two blocks at July 9 to 10 and nowhere else, and the September and October sessions split at October 9. This is how we established the timeline above, which the laboratory log had recorded wrongly. \Cref{sec:tracker} turns the same observation into a detector that runs while the robot works.
